


\pagestyle{DS_FP}
%\NomPrenomNote{}

\bigskip

\exods{\getpts{Exo1}}
%%% POLYNESIE 2025 - J1 - exo1



\section*{PARTIE A}

Un coeur de réseau est constitué d'un maillage dans lequel les routeurs $R_1$ à $R_6$ sont reliés par des liaisons physiques dont le débit en Mbps (Méga-bits par seconde) est indiqué sur la figure suivante.

\begin{figure}[!h]
\begin{center}
\begin{tikzpicture}[
    scale=0.9, transform shape, % Réduit tout à 70%    
    r/.style={draw, circle, minimum size=8mm},
    lan/.style={draw, rectangle, minimum size=6mm},
    % On renomme le style "nuage" pour éviter de boucler avec la forme "cloud"
    nuage/.style={draw, cloud, cloud puffs=15, cloud puff arc=120, aspect=2, minimum width=3cm, minimum height=2cm},
    lbl/.style={midway, fill=white, font=\footnotesize, inner sep=1pt}
]
    % Positionnement des Routeurs
    \node[r] (R1) at (0,0) {$R_1$};
    \node[r] (R2) at (3,2.5) {$R_2$};
    \node[r] (R3) at (3,-2.5) {$R_3$};
    \node[r] (R4) at (9,2.5) {$R_4$};
    \node[r] (R5) at (9,-2.5) {$R_5$};
    \node[r] (R6) at (12,0) {$R_6$};

    % Positionnement des LAN
    \node[lan] (L1) at (-3,0) {$LAN_1$};
    \node[lan] (L2) at (-3,2.5) {$LAN_2$};
    \node[lan] (L3) at (-3,-2.5) {$LAN_3$};
    \node[lan] (L4) at (12,2.5) {$LAN_4$};
    \node[lan] (L5) at (12,-2.5) {$LAN_5$};

    % Internet (on utilise le style 'nuage')
    \node[nuage] (Net) at (15.5,0) {INTERNET};

    % Connexions LAN -> Routeurs
    \draw (L1) -- (R1); \draw (L2) -- (R2); \draw (L3) -- (R3);
    \draw (L4) -- (R4); \draw (L5) -- (R5);
    \draw (R6) -- (Net);

    % Connexions entre Routeurs avec débits
    \draw (R1) -- node[lbl, sloped] {10 Mbps} (R2);
    \draw (R1) -- node[lbl, pos=0.3, sloped] {1 Mbps} (R4);
    \draw (R1) -- node[lbl, pos=0.3, sloped] {1 Mbps} (R5);
    \draw (R2) -- node[lbl] {10 Mbps} (R4);
    \draw (R2) -- node[lbl, pos=0.2] {100 Mbps} (R3);
    \draw (R2) -- node[lbl, pos=0.2, sloped] {10 Mbps} (R5);
    \draw (R3) -- node[lbl, pos=0.2, sloped] {100 Mbps} (R4);
    \draw (R4) -- node[lbl] {100 Mbps} (R5);
    \draw (R5) -- node[lbl, sloped] {100 Mbps} (R6);
\end{tikzpicture}
\end{center}
\caption{Schéma du réseau}\label{exemplesresolutions4}
\end{figure}

Derrière chaque routeur $R_1$ à $R_5$, un switch distribue un réseau local $LAN_1$ à $LAN_5$.
\smallskip
Les protocoles utilisés sont le protocole RIP et le protocole OSPF, qui minimisent respectivement le nombre de routeurs traversés et la somme des coûts des liaisons physiques empruntées.

\medskip

\begin{enumerate}[series=monDS, resume]

\begin{minipage}[h]{0.5\textwidth}
    \item Compléter la table de routage de $R_1$ sachant que le protocole de routage RIP est utilisé. (/\getpts{Q1})

    La destination est le routeur à atteindre, le prochain saut est le premier routeur traversé et la distance est le nombre de routeurs traversés.

\medskip

\item Un utilisateur du réseau local $LAN_1$ interroge, via son navigateur web, un serveur d'Internet. Détailler, suivant le protocole RIP, la route suivie dans le maillage par la requête à destination de ce serveur. (/\getpts{Q2})

\begin{blocvisiblex}[height=1]
{\quad}
 $R_1 \rightarrow R_5 \rightarrow R_6 \rightarrow INTERNET$.
    \end{blocvisiblex}
\end{minipage}\hfill{}
\begin{minipage}[h]{0.44\textwidth}
\begin{blocvisiblex}[height=5]
    {\centering
    \renewcommand{\arraystretch}{1.3}
    \begin{tabular}{|c|c|c|}
        \hline
        \multicolumn{3}{|c|}{\textbf{Table de routage $R_1$}} \\ \hline
        destination & prochain saut & distance \\ \hline
        $R_2$       & $R_2$         & 0        \\ \hline
        $R_3$       & $R_2$         & 1        \\ \hline
        $R_4$       &               &          \\ \hline
        $R_5$       &               &          \\ \hline
        $R_6$       &               &          \\ \hline
    \end{tabular}}
\centering
    \renewcommand{\arraystretch}{1.3}
    \begin{tabular}{|c|c|c|}
        \hline
        \multicolumn{3}{|c|}{\textbf{Table de routage $R_1$}} \\ \hline
        destination & prochain saut & distance \\ \hline
        $R_2$       & $R_2$         & 0        \\ \hline
        $R_3$       & $R_2$         & 1        \\ \hline
        $R_4$       & $R_4$         & 0        \\ \hline
        $R_5$       & $R_5$         & 0        \\ \hline
        $R_6$       & $R_5$         & 1        \\ \hline
        
    \end{tabular}
\end{blocvisiblex}
    
\end{minipage}
 \end{enumerate}

 \medskip



\smallskip

\begin{minipage}{0.50\linewidth}
Le coût d'une liaison est donné par la relation coût $\dfrac{10^2}{\text{débit}}$ où le débit est le débit de la liaison en Mbps.
\end{minipage}
\hfill
\begin{minipage}{0.45\linewidth}
    \centering
    \renewcommand{\arraystretch}{1.3} % Pour l'espacement vertical
    \begin{tabular}{|c|c|c|c|}
        \hline
        \multicolumn{4}{|c|}{\textbf{Coût des liaisons depuis $R_1$}} \\ \hline
        routeur & $R_2$ & $R_4$ & $R_5$ \\ \hline
        coût    & 10    & 100   & 100   \\ \hline
    \end{tabular}
\end{minipage}

\medskip
%%%%%%%%%%%%%%%%%%%%%%%%%%%%%%ù


%%%%%%%%%%%%%%%%%%%%%%%%%%%%%%%%%
\begin{minipage}[h]{0.50\linewidth}
\begin{enumerate}[resume=monDS]
        \item Compléter la table de routage de $R_1$ sachant que le protocole de routage OSPF est utilisé. (/\getpts{Q3})

Un utilisateur du réseau local $LAN_1$ interroge, via son navigateur web, un serveur d'Internet.

\item Donner, selon le protocole OSPF, la route suivie dans le maillage par la requête à destination de ce serveur. (/\getpts{Q4})
    \end{enumerate}
    
\begin{blocvisiblex}[height=1]   
    {}
    $R_1 \rightarrow R_2 \rightarrow R_3 \rightarrow R_4 \rightarrow R_5 \rightarrow R_6 \rightarrow INTERNET$.
    \end{blocvisiblex}
\end{minipage}
\begin{minipage}[h]{0.45\linewidth}
    \begin{blocvisiblex}[height=5]
{    \centering
    \renewcommand{\arraystretch}{1.3} % Pour aérer les cellules comme sur l'image
    \begin{tabular}{|c|c|c|}
        \hline
        \multicolumn{3}{|c|}{\textbf{Table de routage $R_1$}} \\ \hline
        destination & prochain saut & distance \\ \hline
        $R_2$       & $R_2$         & 10        \\ \hline
        $R_3$       &               &          \\ \hline
        $R_4$       &               &          \\ \hline
        $R_5$       &               &          \\ \hline
        $R_6$       &               &          \\ \hline
    \end{tabular}}
\centering
    \renewcommand{\arraystretch}{1.3} % Pour aérer les cellules comme sur l'image
    \begin{tabular}{|c|c|c|}
        \hline
        \multicolumn{3}{|c|}{\textbf{Table de routage $R_1$}} \\ \hline
        destination & prochain saut & distance \\ \hline
        $R_2$       & $R_2$         & 10        \\ \hline
        $R_3$       & $R_2$              &  11        \\ \hline
        $R_4$       & $R_2$              &  12        \\ \hline
        $R_5$       & $R_2$              &  13        \\ \hline
        $R_6$       & $R_2$              &  14        \\ \hline
        
    \end{tabular}
\end{blocvisiblex}
\end{minipage}




\medskip
\pagebreak

    Le protocole de routage OSPF est toujours utilisé et le routeur $R_2$ tombe en panne.

    

    \begin{enumerate}[resume=monDS]
    \item Donner la nouvelle route suivie dans le maillage par une requête partant du réseau $LAN_1$ à destination d'Internet et en donner la nouvelle valeur de la distance.   (/\getpts{Q5})
    \end{enumerate}
    \begin{blocvisiblex}[height=1]
        {}
        $R_1 \rightarrow R_5 \rightarrow R_6 \rightarrow INTERNET$ et la distance est de 101.
    \end{blocvisiblex}
    
\section*{PARTIE B}

Le réseau $LAN_1$ est supposé disposer d'un serveur DHCP (Dynamic Host Control Protocol) gérant l'attribution d'adresses IPv4. 

Il est rappelé, ici, qu'une adresse IPv4 est une séquence de 4 octets, généralement représentée par les 4 entiers correspondants en écriture décimale, séparés par des points (notation décimale pointée). 

Un réseau est caractérisé par des machines dont les adresses ont en commun les $n$ mêmes premiers bits. 

La notation CIDR du réseau a la forme “adresse réseau / $n$”. Appliqué à une adresse du réseau, le masque permet de calculer le préfixe réseau commun. Il est constitué de 4 octets consécutifs, dont (de gauche à droite) les $n$ premiers bits sont égaux à 1 et les suivants à 0. On obtient l'adresse du réseau en effectuant un ET logique, bit à bit, entre chaque octet composant l'adresse d'une machine et l'octet qui lui correspond dans le masque.
 
\medskip

Exemple de \texttt{ET}:

\begin{verbatim}   
    l'entier 192 s'écrit : 1 1 0 0 0 0 0 0
    l'entier 255 s'écrit : 1 1 1 1 1 1 1 1
---------------------------------------------------------
   ET logique, bit à bit : 1 1 0 0 0 0 0 0 ( soit l'entier 192 )
\end{verbatim}
\begin{enumerate}[resume=monDS]
    \item Compléter le tableau suivant correspondant à une machine d'un réseau d'adresse \texttt{192.168.1.100} et de masque \texttt{255.192.0.0}. Ce tableau détermine l'adresse du réseau (en binaire puis en décimale pointée). (/\getpts{Q6})

    \begin{blocvisiblex}[height=4.5]
    {
    
    \centering
    \renewcommand{\arraystretch}{1.5} % Pour l'espacement vertical
    \begin{tabular}{|c|c|c|c|c|}
        \hline
        \multicolumn{5}{|c|}{\textbf{Calcul d'une adresse de réseau}} \\ \hline
        machine (binaire)        & 11000000 & 10101000 &          &          \\ \hline
        masque (binaire)         & 11111111 & 11000000 & 00000000 & 00000000 \\ \hline
        réseau (binaire)         & 11000000 &          &          &          \\ \hline
        réseau (déc.\ pointée)    & 192      &          &          &          \\ \hline
    \end{tabular}
}
\centering
    \renewcommand{\arraystretch}{1.5} % Pour l'espacement vertical
    \begin{tabular}{|c|c|c|c|c|}
        \hline
        \multicolumn{5}{|c|}{\textbf{Calcul d'une adresse de réseau}} \\ \hline
        machine (binaire)        & 11000000 & 10101000 & 00000001 & 01100100\\ \hline
        masque (binaire)         & 11111111 & 11000000 & 00000000 & 00000000 \\ \hline
        réseau (binaire)         & 11000000 & 10000000 & 00000000 & 00000000 \\ \hline
        réseau (déc.\ pointée)    & 192      & 128      & 0        & 0        \\ \hline
    \end{tabular}
    \end{blocvisiblex}

\end{enumerate}

On obtient le complémentaire d'un octet en échangeant respectivement chacun des 0 et des 1 qui le composent, par un 1 ou un 0.

Par exemple, le complémentaire de \verb|1 1 0 0 1 0 1 0| est \verb|0 0 1 1 0 1 0 1|.

Par un procédé analogue à celui de la question précédente, l'adresse de broadcast d'un réseau est obtenue en effectuant un OU logique bit à bit entre chaque octet composant l'adresse réseau et le complémentaire de l'octet correspondant dans l'écriture binaire du masque.

\medskip

Exemple de \texttt{OU}:
\begin{verbatim}   
           12 s'écrit : 0 0 0 0 1 1 0 0
            9 s'écrit : 0 0 0 0 1 0 0 1
------------------------------------------------------------
OU logique, bit à bit : 0 0 0 0 1 1 0 1 ( soit l'entier 13 )
\end{verbatim}

\pagebreak

\begin{enumerate}[resume=monDS]
    \item Compléter le tableau ci-après pour déterminer l'adresse de broadcast du réseau suivant (en binaire puis en décimale pointée). (/\getpts{Q7})

    \begin{blocvisiblex}[height=5]
    {\centering
    \renewcommand{\arraystretch}{1.5} % Ajustement de la hauteur des lignes
    \begin{tabular}{|c|c|c|c|c|}
        \hline
        \multicolumn{5}{|c|}{\textbf{Calcul d’une adresse de broadcast}} \\ \hline
        réseau (binaire) & 11000000 & 10000000 & 00000000 & 00000000 \\ \hline
        masque (binaire) & 11111111 & 11000000 & 00000000 & 00000000 \\ \hline
        complément du masque (binaire) & & & & \\ \hline
        broadcast (binaire) & & & & \\ \hline
        broadcast (déc.\ pointée) & & & & \\ \hline
    \end{tabular}}
    \centering
    \renewcommand{\arraystretch}{1.5} % Ajustement de la hauteur des lignes
    \begin{tabular}{|c|c|c|c|c|}
        \hline
        \multicolumn{5}{|c|}{\textbf{Calcul d’une adresse de broadcast}} \\ \hline
        réseau (binaire) & 11000000 & 10000000 & 00000000 & 00000000 \\ \hline
        masque (binaire) & 11111111 & 11000000 & 00000000 & 00000000 \\ \hline
        complément du masque (binaire) & 00000000 & 00111111 & 11111111 & 11111111 \\ \hline
        broadcast (binaire) & 11000000 & 10111111 & 11111111 & 11111111 \\ \hline
        broadcast (déc. pointée) & 192 & 191 & 255 & 255 \\ \hline
        
    \end{tabular}
        
    \end{blocvisiblex}
    


\end{enumerate}

\medskip

Une machine du réseau $LAN_1$ a reçu du serveur DHCP l'adresse \verb|172.16.1.100|, avec le masque \verb|255.255.0.0|.

Sa passerelle par défaut a pour adresse \verb|172.16.255.254|.

\begin{enumerate}[resume=monDS]
    \item En déduire les informations suivantes: (/\getpts{Q8})
    \begin{itemize}
        \item l'adresse du réseau $LAN_1$ (en décimale pointée);
        
        \begin{blocvisiblex}[height=1]
            {\quad}
            L'adresse du réseau $LAN_1$ est \texttt{172.16.0.0}.
        \end{blocvisiblex}
        \item l'adresse de broadcast (en décimale pointée);
        \begin{blocvisiblex}[height=1]
            {\quad}
            L'adresse de broadcast du réseau $LAN_1$ est \texttt{172.16.255.255}.
        \end{blocvisiblex}
        
        \item le nombre total d'adresses pouvant être distribuées par le serveur, en ne tenant pas compte des éventuelles restrictions possiblement posées par l'administrateur du réseau.
        \begin{blocvisiblex}[height=1.5]
            {\quad}
            Le nombre total d'adresses pouvant être distribuées par le serveur DHCP est de $2^{16} - 2 = 65534$ adresses, car il faut exclure l'adresse du réseau et l'adresse de broadcast.
        \end{blocvisiblex}
    \end{itemize}
\end{enumerate}

\bigskip

On souhaite, désormais, simuler en Python le fonctionnement du réseau $LAN_1$.

On donne, ci-après, les documentations des méthodes \verb|split| et \verb|join|.

\begin{itemize}
    \item la documentation de la méthode \verb|split| de l'objet \verb|str| est la suivante:
\begin{CodePythontexAlt}[TaillePolice=\large, Largeur=1\linewidth,Centre,Lignes=true,PremLigne=1]{}
    split(self, séparateur) 
        Renvoie la liste des sous-chaines de caractères délimitées par 
        le séparateur dans l'instance courante de chaîne de caractères. 
        >>> 'ab-pq-rs'.split('-') 
        ['ab', 'pq', 'rs']
    \end{CodePythontexAlt}

    \medskip

    \item la documentation de la méthode \verb|join| de l'objet \verb|str|:
\begin{CodePythontexAlt}[TaillePolice=\large, Largeur=1\linewidth,Centre,Lignes=true,PremLigne=1]{}
    join(self, iterable) 
        Fusionne les chaînes de caractères contenues dans iterable en 
        le liant par la sous-chaîne depuis laquelle cette méthode est appelée. 
        >>>'.'.join(['ab', 'pq', 'rs']) 
        'ab.pq.rs'
    \end{CodePythontexAlt}

\end{itemize}

On commence par créer la classe IPv4 suivante.

\begin{CodePythontexAlt}[TaillePolice=\large, Largeur=1\linewidth,Centre,Lignes=true,PremLigne=1]{}
class IPv4(object):
    def __init__(self, adresse:str):
        """ Constructeur de la classe de calcul sur une adresse IPv4 
        dont la notation décimale pointée est passée en paramètre. """
        self.adresse = adresse
    def octets(self)->list[int]:
        """ Découpe l'adresse décimale pointée en la liste des 4 entiers 
        correspondants. 
        >>> ip01 = IPv4('192.168.1.100')
        >>> ip01.octets()
        [192, 168, 1, 100]        """
        return [int(i) for i in self.adresse.split(".")]

   \end{CodePythontexAlt}

\medskip

   Pour mémoire, l'opérateur \texttt{\&} entre deux entiers effectue le ET logique bit à bit de la représentation binaire de ces entiers et renvoie l'entier correspondant à la représentation binaire ainsi obtenue. Exemple:
   
   \begin{verbatim}
    >>> 192 & 255 
    192
   \end{verbatim}

   \begin{enumerate}[resume=monDS]
    \item Compléter les lignes 11 et 12 du code de la méthode \verb|masquer| de la classe \verb|IPv4|, donné ci-dessous. La méthode doit correspondre au docstring. (/\getpts{Q9})
    
\begin{CodePythontexAlt}[TaillePolice=\large, Largeur=1\linewidth,Centre,Lignes=true,PremLigne=1]{}
    def masquer(self, masque: str)->str:
        """ Détermine le préfixe masqué de l'adresse, 
        le masque (décimal pointé) étant passé en paramètre. 
        >>> ip01 = IPv4('192.168.1.100')
        >>> ip01.masquer('255.192.0.0')
        '192.128.0.0'        """
        tmp = []
        ip = self.octets()
        crible = IPv4(masque).octets()
        for i in range(4):
            tmp.append(...... & ......)
        return ".".join(.........)
   \end{CodePythontexAlt}

   \begin{blocvisible}[height=1.5]
        \begin{lstlisting}
    Ligne 11: tmp.append(ip[i] & crible[i])
    Ligne 12: ".".join([str(k) for k in tmp])        
    \end{lstlisting}
    \end{blocvisible}


        \begin{CodePythontexAlt}[TaillePolice=\large, Largeur=1\linewidth,Centre,Lignes=true,PremLigne=1]{}
    def adresse_suivante(self, adresse_max:str)->str:
        """ Détermine l'adresse décimale pointée suivant immédiatement 
        l'adresse courante, sous réserve d'existence d'une adresse disponible 
        >>> add = IPv4('192.168.1.100')
        >>> add.adresse_suivante('192.168.1.254')
        '192.168.1.101'
        >>> add = IPv4('192.168.1.255')
        >>> add.adresse_suivante('192.168.255.254')
        '192.168.2.0'        """
        assert self.adresse < adresse_max
        liste_courante = self.octets()
        liste_suivante = list()
        retenue = 1
        for index in range(4):
            somme = liste_courante[3 - index] + retenue
            valeur, retenue = ......, ......
            liste_suivante =......
        return '.'.join(liste_suivante)
   \end{CodePythontexAlt}
        
    

   \medskip

\item Compléter les lignes 16 et 17 du code de la méthode \verb|adresse_suivante| de la classe \verb|IPv4|, donné ci-dessous. La méthode doit correspondre au docstring. (/\getpts{Q10})

  \begin{blocvisible}[height=1.5]
        \begin{lstlisting}
    Ligne 16: valeur, retenue = somme%256, somme//256
    Ligne 17: liste_suivante = [str(valeur)] + liste_suivante     
    \end{lstlisting}
    \end{blocvisible}


   \end{enumerate}
